%% kapa: exam=1 round=35 year=2024 period=2 form=A subject=appraisal_law subject_ko=감정평가관계법규 seq=8 number=8
%% source: Q-Net 감정평가사 기출문제 (cst00309) artlSeq=5213229 file=raw/1008/5213229/2교시 문제지 원본.pdf
%% conversion: kordoc 4.24.51 -> tools/gen1.py (auto); review=proofread-sonnet; flags=none
\begin{question}{8}
국토의 계획 및 이용에 관한 법령상 개발행위허가를 받은 자가 행정청인 경우 개발행위에 따른 공공시설의 귀속에 관한 설명으로 옳지 않은 것은?
\begin{choices}
  \item 개발행위허가를 받은 자가 새로 설치한 공공시설은 그 시설을 관리할 관리청에 무상으로 귀속된다.
  \item 개발행위허가를 받은 자가 기존의 공공시설에 대체되는 공공시설을 설치한 경우 종래의 공공시설은 개발행위허가를 받은 자에게 무상으로 귀속된다.
  \item 새로 설치된 공공시설의 귀속시점은 준공검사를 받은 날이다.
  \item 개발행위허가를 받은 자는 개발행위가 끝나 준공검사를 마친 때에는 해당 시설의 관리청에 공공시설의 종류와 토지의 세목을 통지하여야 한다.
  \item 개발행위허가를 받은 자는 그에게 귀속된 공공시설의 처분으로 인한 수익금을 도시ㆍ군계획사업 외의 목적에 사용하여서는 아니 된다.
\end{choices}
\end{question}
\answer{3}
